\documentclass[11pt]{article}

\usepackage[margin=1in]{geometry}
\usepackage{amsmath,amssymb,amsthm}
\usepackage{algorithm}
\usepackage{algpseudocode}
\usepackage{enumitem}
\usepackage{hyperref}
\usepackage{xcolor}
\usepackage{tikz}
\usepackage{booktabs}

\hypersetup{colorlinks=true, linkcolor=blue, urlcolor=blue, citecolor=blue}

\newtheorem{theorem}{Theorem}[section]
\newtheorem{lemma}[theorem]{Lemma}
\newtheorem{proposition}[theorem]{Proposition}
\newtheorem{corollary}[theorem]{Corollary}
\newtheorem{definition}[theorem]{Definition}
\newtheorem{example}[theorem]{Example}

\title{TOAE209/TOAH209 -- Design and Analysis of Algorithms\\[4pt]
\large Week 4: Greedy Algorithms I}
\author{Foreign Trade University}
\date{}

\begin{document}
\maketitle
\tableofcontents

\section{Introduction: The Greedy Paradigm}

This week begins our study of \textbf{greedy algorithms}: algorithms that build up a
solution piece by piece, at each step making the choice that looks best \emph{right now},
according to some simple local rule, and never reconsidering that choice later.

Greedy algorithms are appealing because they are usually short, fast, and easy to
implement -- often just "sort the input by some key, then scan once." The hard part is
never the implementation; it is the \textbf{proof}. For almost any greedy rule one can
imagine, there is a plausible-sounding argument for why it should work, and (often) a
counter-example showing it does not. Distinguishing the two requires a careful
correctness proof.

\subsection{Two proof techniques}

This course (following Kleinberg \& Tardos, Ch.\ 4) uses two main templates for proving a
greedy algorithm correct.

\begin{itemize}
    \item \textbf{Greedy stays ahead.} Define a natural measure of progress (e.g., "after
    $k$ steps, the greedy solution has done at least as well as any other algorithm's
    first $k$ steps, according to some measure"). Prove by induction on $k$ that the
    greedy solution stays "ahead of" or "tied with" every other feasible solution at every
    step. Conclude that, at the end, greedy is at least as good as any alternative.

    \item \textbf{Exchange argument.} Take an arbitrary optimal solution $O$ that may
    differ from the greedy solution $G$. Show that $O$ can be transformed into $G$ by a
    sequence of local "exchanges" -- swaps that do not make the solution any worse (and
    sometimes strictly improve it, or at least do not violate feasibility). Since each
    exchange preserves (or improves) optimality, and the sequence of exchanges turns $O$
    into $G$, $G$ must also be optimal.
\end{itemize}

Both techniques appear repeatedly in this course. Today we will see the exchange argument
in its cleanest forms, applied to interval scheduling and to minimizing lateness.

\subsection{Why naive greedy rules often fail}

A recurring theme today is that \emph{many plausible greedy rules are wrong}. For a
problem with several "obvious" greedy orderings (by size, by deadline, by start time, by
value, ...), usually only one (or a specific combination) is provably correct, and the
others fail on small counter-examples. Part of learning to design greedy algorithms is
learning to quickly construct such counter-examples -- which is exactly what several of
this week's exercises ask you to do.

\section{Interval Scheduling}

\subsection{Problem statement}

\begin{definition}[Interval Scheduling Problem]
We are given a set of $n$ \emph{requests} (or \emph{jobs}) $1, 2, \ldots, n$, where request
$i$ has a \emph{start time} $s_i$ and a \emph{finish time} $f_i$, with $s_i < f_i$. Two
requests $i$ and $j$ are \emph{compatible} if the intervals $[s_i, f_i)$ and $[s_j, f_j)$
do not overlap (i.e., $s_i \ge f_j$ or $s_j \ge f_i$). The goal is to select a
\textbf{maximum-size} subset of mutually compatible requests.
\end{definition}

Think of each request as a meeting that needs a single shared room: we want to schedule as
many meetings as possible in that one room.

\subsection{Why several natural greedy rules fail}

Before presenting the correct algorithm, it is worth listing some tempting rules and why
each one can fail:

\begin{itemize}
    \item \textbf{Earliest start time first.} Counter-example: one request starting very
    early but lasting all day blocks out many short requests that could otherwise all fit.
    \item \textbf{Shortest interval first.} A short interval can overlap (and thus
    eliminate) two or more longer, mutually compatible intervals -- Problem 6 in this
    week's exercises asks you to construct exactly such an instance.
    \item \textbf{Fewest conflicts first.} Plausible, but the bookkeeping is complex and
    (perhaps surprisingly) it is not needed: a much simpler rule turns out to be optimal.
\end{itemize}

\subsection{The earliest-finish-time-first algorithm}

\begin{algorithm}[h]
\caption{Earliest-Finish-Time-First (Interval Scheduling)}
\begin{algorithmic}[1]
\State Sort requests by finish time: $f_1 \le f_2 \le \cdots \le f_n$
\State $A \gets \emptyset$
\For{$i = 1$ to $n$}
    \If{request $i$ is compatible with every request currently in $A$}
        \State $A \gets A \cup \{i\}$
    \EndIf
\EndFor
\State \Return $A$
\end{algorithmic}
\end{algorithm}

Because requests are processed in order of finish time, "compatible with everything in
$A$" simplifies to "starts no earlier than the finish time of the most recently added
request" -- so the inner check is $O(1)$ given the previous selection, and the whole
algorithm runs in $O(n \log n)$ (dominated by the sort). Problem 3 asks you to implement
exactly this.

\subsection{Correctness: an exchange argument}

\begin{theorem}
The earliest-finish-time-first algorithm produces a maximum-size set of mutually
compatible requests.
\end{theorem}

\begin{proof}
Let $i_1, i_2, \ldots, i_k$ be the requests selected by the greedy algorithm, in the order
selected (so $f_{i_1} < f_{i_2} < \cdots < f_{i_k}$, and $f_{i_1}$ is the smallest finish
time among \emph{all} requests). Let $j_1, j_2, \ldots, j_m$ be the requests of an
arbitrary optimal solution, also sorted by finish time. We show $k \ge m$, which (since
the greedy solution is always feasible) implies $k = m$.

\textbf{Claim (greedy stays ahead).} For every $r$ with $1 \le r \le \min(k,m)$,
$f_{i_r} \le f_{j_r}$.

\textit{Proof of claim, by induction on $r$.}
\begin{itemize}
    \item \textbf{Base case ($r=1$).} $i_1$ was chosen because it has the globally
    smallest finish time among \emph{all} $n$ requests. In particular $f_{i_1} \le
    f_{j_1}$.
    \item \textbf{Inductive step.} Suppose $f_{i_{r-1}} \le f_{j_{r-1}}$. Since $j_r$ is
    compatible with $j_{r-1}$ in the optimal solution, $s_{j_r} \ge f_{j_{r-1}} \ge
    f_{i_{r-1}}$. So $j_r$ is compatible with $i_{r-1}$ -- and hence also compatible with
    every $i_1, \ldots, i_{r-1}$ (which all finish no later than $f_{i_{r-1}}$).
    Therefore, when the greedy algorithm considers requests in increasing order of finish
    time, $j_r$ (or some other request with finish time $\le f_{j_r}$) was a
    \emph{feasible} candidate to extend $i_1, \ldots, i_{r-1}$ at the point where the
    greedy algorithm chose $i_r$. Since greedy picks the smallest-finish-time compatible
    request available, $f_{i_r} \le f_{j_r}$.
\end{itemize}

\textbf{Finishing the proof.} Suppose for contradiction $k < m$. Apply the claim with $r =
k$: $f_{i_k} \le f_{j_k}$. Since $j_{k+1}$ exists (because $m > k$) and is compatible with
$j_k$, we have $s_{j_{k+1}} \ge f_{j_k} \ge f_{i_k}$, so $j_{k+1}$ is compatible with
$i_1,\ldots,i_k$ -- meaning the greedy algorithm \emph{could} have added $j_{k+1}$ (or
some request with an even smaller finish time) after $i_k$, contradicting the fact that
greedy's selection ended at $i_k$ (greedy only stops adding a request when no compatible
one remains, but $j_{k+1}$ was compatible). Hence $k \ge m$, and since $A$ is feasible,
$k \le m$ as well (an optimal solution has size $\ge$ any feasible solution), so $k = m$.
\end{proof}

Problems 4--5 ask you to verify this empirically: implement a brute-force optimum for
small instances and confirm the greedy result always has the same size.

\section{Interval Partitioning}

\subsection{Problem statement}

\begin{definition}[Interval Partitioning Problem]
Given $n$ requests (intervals) as before, partition them into the \emph{minimum} number of
\emph{resources} (e.g.\ classrooms) such that no two requests assigned to the same
resource overlap. Equivalently: find the minimum number of "colors" needed to color the
intervals so that overlapping intervals get different colors.
\end{definition}

This differs from Interval Scheduling: here \emph{every} request must be scheduled
(on some resource), but we get to use as many resources as needed -- we just want to
minimize that number.

\subsection{A lower bound: depth}

\begin{definition}[Depth]
The \emph{depth} of a set of intervals is the maximum, over all points in time $t$, of the
number of intervals containing $t$.
\end{definition}

\begin{lemma}
\label{lem:depth-lower-bound}
No assignment of intervals to resources (with no overlaps on a single resource) can use
fewer than $\mathrm{depth}$ resources.
\end{lemma}

\begin{proof}
If $d = \mathrm{depth}$, then there is some time $t$ contained in $d$ distinct intervals.
These $d$ intervals are pairwise overlapping (they all contain $t$), so no two of them can
share a resource. Hence at least $d$ distinct resources are required.
\end{proof}

\subsection{The earliest-start-time-first algorithm}

\begin{algorithm}[h]
\caption{Earliest-Start-Time-First (Interval Partitioning)}
\begin{algorithmic}[1]
\State Sort intervals by start time: $s_1 \le s_2 \le \cdots \le s_n$
\State $d \gets 0$ \Comment{number of resources allocated so far}
\State $H \gets \emptyset$ \Comment{min-heap of (end\_time, resource\_id) for "busy" resources}
\For{$i = 1$ to $n$}
    \If{$H \ne \emptyset$ and $\min(H).\mathrm{end\_time} \le s_i$}
        \State Pop the resource $r$ with smallest end\_time from $H$
        \State Assign interval $i$ to resource $r$
    \Else
        \State $d \gets d + 1$; assign interval $i$ to a new resource $d$
    \EndIf
    \State Push $(f_i, \text{resource of } i)$ onto $H$
\EndFor
\State \Return $d$ and the assignment
\end{algorithmic}
\end{algorithm}

Using a binary min-heap keyed by end time, each interval requires $O(\log n)$ heap
operations, so the whole algorithm runs in $O(n \log n)$. Problem 7 asks you to implement
this directly.

\subsection{Correctness: the greedy uses exactly \texorpdfstring{$\mathrm{depth}$}{depth} resources}

\begin{theorem}
The earliest-start-time-first algorithm uses exactly $d = \mathrm{depth}$ resources, and
this is optimal.
\end{theorem}

\begin{proof}
By Lemma~\ref{lem:depth-lower-bound}, no algorithm can use fewer than $\mathrm{depth}$
resources, so it suffices to show the greedy algorithm uses \emph{at most}
$\mathrm{depth}$ resources.

Suppose the greedy algorithm allocates a new resource (resource number $d$) when
processing interval $i$. This happens only when, at the moment interval $i$ is considered,
\emph{every} resource $1, \ldots, d-1$ currently holds an interval that has not yet
finished, i.e.\ for each resource $r < d$ there is some interval $j_r$ with $s_{j_r} \le
s_i < f_{j_r}$ (it started no later than $s_i$ and finishes after $s_i$). Together with
interval $i$ itself (which also "contains" $s_i$, since $s_i < f_i$), this gives $d$
intervals -- the $d-1$ intervals $j_1, \ldots, j_{d-1}$ plus interval $i$ -- all containing
the point $s_i$. Hence $\mathrm{depth} \ge d$.

So whenever the greedy algorithm uses a resource numbered $d$, we have $d \le
\mathrm{depth}$. The total number of resources used is therefore at most
$\mathrm{depth}$, and combined with the lower bound, it equals $\mathrm{depth}$ exactly.
\end{proof}

Problem 8 asks you to verify this equality empirically on random instances, and Problems
9--10 explore related feasibility questions (can $k$ resources suffice?) and stress-test
the implementation.

\section{Scheduling to Minimize Lateness}

\subsection{Problem statement}

\begin{definition}[Minimizing Lateness Problem]
We are given $n$ jobs, all available for processing at time $0$. Job $i$ has a
\emph{processing time} (length) $t_i > 0$ and a \emph{deadline} $d_i$. A single resource
processes jobs one at a time, without preemption (once started, a job runs to
completion). For a schedule, let $f_i$ be the finish time of job $i$. The
\emph{lateness} of job $i$ is $\ell_i = \max(0, f_i - d_i)$. The goal is to find a
schedule minimizing $L = \max_i \ell_i$, the \textbf{maximum lateness}.
\end{definition}

\subsection{Why naive rules fail here too}

\begin{itemize}
    \item \textbf{Shortest job first (ignoring deadlines).} A very short job with a very
    tight deadline can be starved by a slightly longer job with no urgency, if the longer
    job is scheduled first merely because... no wait, shortest-first would schedule the
    short job first in that case. The real failure mode: a job with a generous deadline
    but very short length gets scheduled first, "wasting" the early time slots that a job
    with a tight deadline desperately needed. Problem 14 asks you to construct such an
    instance.
    \item \textbf{Smallest slack first (i.e., $d_i - t_i$ smallest first).} This rule is
    actually used for some related problems but can be shown to fail here too in certain
    edge cases -- we will not pursue it further in this course.
\end{itemize}

\subsection{The earliest-deadline-first algorithm}

\begin{algorithm}[h]
\caption{Earliest-Deadline-First (Minimizing Lateness)}
\begin{algorithmic}[1]
\State Sort jobs by deadline: $d_1 \le d_2 \le \cdots \le d_n$
\State $t \gets 0$
\For{$i = 1$ to $n$}
    \State Schedule job $i$ to run in $[t, t + t_i)$
    \State $f_i \gets t + t_i$
    \State $t \gets f_i$
\EndFor
\State \Return the schedule and $L = \max_i \max(0, f_i - d_i)$
\end{algorithmic}
\end{algorithm}

This is exactly "sort by deadline, then run them back to back" -- $O(n \log n)$ total.
Problem 11 asks you to implement this.

\subsection{Two structural facts about optimal schedules}

\begin{definition}[Idle time and inversions]
A schedule has \emph{idle time} if the resource is idle at some time $t$ while some job
has not yet been completed (and could have started by time $t$). A schedule (given as an
ordering of jobs) has an \emph{inversion} if there exist jobs $i, j$ with $i$ scheduled
before $j$ but $d_i > d_j$.
\end{definition}

\begin{lemma}
\label{lem:no-idle}
There exists an optimal schedule with no idle time and, among schedules with no idle
time, an optimal one exists with no inversions.
\end{lemma}

\begin{proof}[Proof sketch]
\textbf{No idle time.} If a schedule has idle time, we can shift every job after the idle
period earlier (removing the gap) without increasing any finish time, hence without
increasing any lateness. So removing idle time never makes the schedule worse.

\textbf{No inversions (exchange argument).} Suppose an idle-free schedule has an
inversion: some adjacent pair of jobs $i$ (scheduled right before $j$) with $d_i > d_j$
(any inversion must contain an adjacent inverted pair, by a standard argument -- if $i$
before $j$ is an inversion but they are not adjacent, there is some adjacent inverted pair
between them). Swap $i$ and $j$. All other jobs' finish times are unchanged. Job $j$ now
finishes earlier than before (it moved earlier in the order), so its lateness can only
decrease. Job $i$ now finishes at the time $j$ used to finish, call it $f$. We claim job
$i$'s new lateness $\max(0, f - d_i)$ is \emph{at most} the larger of the two jobs' old
latenesses. Concretely, before the swap, job $j$ (the later one) finished at time $f$, so
its old lateness was $\max(0, f - d_j)$. Since $d_i > d_j$, we have $f - d_i < f - d_j$, so
$\max(0,f-d_i) \le \max(0, f-d_j)$ -- job $i$'s new lateness is no worse than job $j$'s old
lateness. Hence the maximum lateness over the whole schedule does not increase after the
swap. Repeating this for every inversion (each swap strictly reduces the inversion count)
eventually yields an inversion-free schedule that is no worse than the original.
\end{proof}

\begin{theorem}
The earliest-deadline-first algorithm produces a schedule with the minimum possible
maximum lateness.
\end{theorem}

\begin{proof}
By Lemma~\ref{lem:no-idle}, some optimal schedule $O$ has no idle time and no inversions.
A schedule with no idle time and no inversions must process jobs in non-decreasing order
of deadline (any other order would contain an inversion) -- which is \emph{exactly} the
order produced by earliest-deadline-first. Since all idle-free schedules with the same job
order have identical finish times (jobs run back-to-back starting at time 0), $O$ and the
earliest-deadline-first schedule are the same schedule (up to tie-breaking among jobs with
equal deadlines, which does not affect $L$). Hence earliest-deadline-first achieves the
optimal $L$.
\end{proof}

Problems 12--13 ask you to verify this against a brute-force search over all $n!$
orderings for small $n$, and Problem 15 asks you to directly measure inversions in EDF's
output (it should always be zero) and in perturbed orderings.

\section{Optimal Caching}

\subsection{Problem statement}

\begin{definition}[Optimal Caching Problem]
We are given a cache that can hold up to $k$ items, and a sequence of $m$ requests $\sigma
= \sigma_1, \sigma_2, \ldots, \sigma_m$, each requesting some item. If a requested item is
already in the cache, this is a \emph{hit} (free); otherwise it is a \emph{miss}, and if
the cache is full, some item currently in the cache must be \emph{evicted} to make room.
The goal (in the \emph{offline} version, where the entire sequence $\sigma$ is known in
advance) is to choose an eviction policy minimizing the total number of misses.
\end{definition}

This models, e.g., a CPU cache, a web browser's page cache, or a database buffer pool:
$k$ is the cache capacity, and each "miss" represents an expensive operation (a memory
fetch, a disk read, a network request) that a hit avoids.

\subsection{The furthest-in-future algorithm}

\begin{algorithm}[h]
\caption{Furthest-In-Future (Belady's MIN / LFD)}
\begin{algorithmic}[1]
\For{each request $\sigma_i$, in order}
    \If{$\sigma_i$ is already in the cache}
        \State (hit; do nothing)
    \Else
        \If{the cache is full}
            \State Evict the item in the cache whose \emph{next} occurrence in
            $\sigma_{i+1}, \sigma_{i+2}, \ldots$ is furthest in the future (or
            that never occurs again)
        \EndIf
        \State Load $\sigma_i$ into the cache (miss)
    \EndIf
\EndFor
\end{algorithmic}
\end{algorithm}

\begin{theorem}[Belady, 1966]
Furthest-in-future is an optimal offline caching policy: it achieves the minimum possible
number of misses on any request sequence.
\end{theorem}

\begin{proof}[Proof sketch -- exchange argument]
Let $\mathrm{FF}$ denote the furthest-in-future schedule of evictions, and let $O$ be any
optimal eviction schedule. We show $O$ can be transformed, step by step, into a schedule
that behaves identically to $\mathrm{FF}$, without increasing the number of misses, via
the following idea: consider the first point at which $\mathrm{FF}$ and $O$ make a
different eviction decision on a miss. At that point, $\mathrm{FF}$ evicts the item $x$
used furthest in the future; suppose $O$ evicts some other item $y$ (so $y$'s next use is
sooner than $x$'s). We can modify $O$ to evict $x$ instead of $y$ at this step, and show
that $O$'s subsequent behavior can be adjusted (by "pretending" to still have $y$ in cache
until it is needed, then evicting whatever $\mathrm{FF}$ would have evicted at that later
point) so that the total number of misses does not increase. This is the technique called
"reduction to a canonical form" -- repeating this argument shows some optimal schedule
agrees with $\mathrm{FF}$ at every step, hence $\mathrm{FF}$ itself is optimal. The full
details are in Kleinberg \& Tardos, Section 4.4.
\end{proof}

\subsection{Online caching: LRU}

Furthest-in-future requires knowing the \emph{future} request sequence -- it is an
\emph{offline} algorithm. In practice, caches must make decisions \emph{online}, without
knowing future requests. The most common online heuristic is \textbf{Least-Recently-Used
(LRU)}: evict the item whose most recent access is furthest in the \emph{past}, on the
heuristic assumption that recently-used items are likely to be used again soon
("temporal locality").

LRU can be far from optimal: Problem 18 asks you to construct a request sequence where LRU
incurs strictly more misses than furthest-in-future, by exploiting a cyclic access pattern
that defeats the "recently used = likely reused" heuristic. (In the worst case, LRU can be
forced to miss on \emph{every single request} in a cyclic pattern over $k+1$ items with a
cache of size $k$, while furthest-in-future does strictly better.)

\section{A Different Flavor: Greedy on Coin Change and Job Sequencing}

Not every greedy algorithm in this chapter is about intervals. Two more classical
examples, explored in Problems 20--22:

\subsection{Coin change}

Given coin denominations (e.g.\ $\{1, 2, 5, 10, 20, 50, 100, 200, 500\}$) and a target
amount, the \textbf{greedy} algorithm repeatedly takes the largest denomination that does
not exceed the remaining amount. For \emph{canonical} coin systems (which include all
standard currency denominations in common use), this greedy algorithm uses the minimum
possible number of coins. However, this is \emph{not} true for arbitrary denomination
sets: with denominations $\{1, 3, 4\}$, greedy makes change for $6$ as $4+1+1$ (3 coins),
while $3+3$ (2 coins) is strictly better. Whether a coin system is canonical is, in
general, a non-trivial property -- proving it requires checking a finite (but possibly
large) set of "witness" amounts.

\subsection{Job sequencing with deadlines and profits}

Each job $i$ has a deadline $d_i$ and a profit $p_i$, takes one unit of time, and earns
$p_i$ only if scheduled at or before $d_i$ on a single machine (with time slots $1, 2, 3,
\ldots$). The greedy algorithm processes jobs in \emph{decreasing order of profit}, and
places each job in the \emph{latest} available slot at or before its deadline (if any slot
is available). The intuition: placing a job as late as possible "wastes" the fewest early
slots, leaving them open for other jobs with tighter deadlines. This greedy algorithm is
optimal, and can be proven correct with another exchange argument (swapping a job into a
later free slot never decreases total profit, since both jobs still meet their
deadlines).

\section{Summary and Looking Ahead}

This week introduced the \textbf{greedy algorithm design paradigm} and two proof
techniques -- \emph{greedy stays ahead} and the \emph{exchange argument} -- through four
case studies:

\begin{itemize}
    \item \textbf{Interval Scheduling}: earliest-finish-time-first is optimal (greedy
    stays ahead).
    \item \textbf{Interval Partitioning}: earliest-start-time-first uses exactly
    $\mathrm{depth}$ resources, which is optimal.
    \item \textbf{Minimizing Lateness}: earliest-deadline-first is optimal (exchange
    argument via inversions).
    \item \textbf{Optimal Caching}: furthest-in-future is optimal offline (exchange
    argument / reduction to canonical form); LRU is a practical but suboptimal online
    heuristic.
\end{itemize}

We also saw, repeatedly, that \emph{not every natural-sounding greedy rule is correct} --
shortest-interval-first and shortest-job-first both fail, and the non-canonical coin
system $\{1,3,4\}$ shows that even "obviously correct" greedy rules (like coin change) can
fail outside their intended domain. Constructing counter-examples is as important a skill
as designing the correct algorithm.

Next week (Greedy II, Kleinberg \& Tardos Ch.\ 4 continued) we apply the same exchange-
argument toolkit to two fundamental graph algorithms: Dijkstra's shortest-path algorithm
and minimum spanning trees (Kruskal's and Prim's algorithms).

\end{document}
